\chapter{Implementación numérica}

En los capítulos anteriores se han presentado las dos metodologías que se van a comparar: por un lado, el modelo Heston--SLV, que modifica la dinámica del subyacente mediante una función de apalancamiento, y por otro lado la metodología de Entropía Relativa Mínima, que modifica la medida de probabilidad asignando nuevos pesos a trayectorias previamente simuladas. En este capítulo se describe cómo se han implementado numéricamente ambos enfoques.\medskip

El objetivo no es todavía analizar los resultados obtenidos, sino fijar con precisión el procedimiento computacional utilizado. La comparación de precios, los errores de ajuste y el comportamiento de ambos métodos frente a opciones \textit{forward-start} se discutirán posteriormente en el capítulo de resultados numéricos.

\section{Datos de entrada y configuración común}

La implementación parte de una superficie inicial de volatilidad implícita asociada a opciones europeas vainilla. En los experimentos considerados se trabaja con un subyacente inicial
\[
S_0=100,
\]
y con tipo de interés constante \(r\). La superficie se representa mediante una malla discreta de strikes y vencimientos,
\[
\{K_1,\ldots,K_M\},
\qquad
\{T_1,\ldots,T_L\},
\]
junto con los valores de volatilidad implícita observados en cada punto,
\[
\widehat{\sigma}(T_\ell,K_m).
\]

A partir de esta superficie se construyen los precios de referencia de las opciones vainilla mediante la fórmula de Black--Scholes. Estos precios actúan como precios de mercado o precios benchmark, dependiendo del experimento considerado. De este modo, ambos métodos se comparan frente a la misma información inicial.\medskip

En el código, las opciones utilizadas son de tipo call y los vencimientos se expresan en años. La configuración puede modificarse para escoger distintos strikes de valoración, distintos vencimientos, número de trayectorias o semilla de simulación. Es importante distinguir entre los instrumentos utilizados para calibrar cada método y los instrumentos utilizados después para valorar. En particular, una vez generadas las trayectorias, es posible valorar opciones con strikes que no hayan formado parte de la calibración.

\section{Implementación del modelo Heston--SLV}

\subsection{Calibración de Heston y construcción de la volatilidad local}

El primer paso en la implementación SLV consiste en calibrar el modelo de Heston a la superficie inicial. La calibración se realiza sobre los parámetros
\[
\kappa,\qquad \theta,\qquad \gamma,\qquad \rho_{S,v},\qquad v_0,
\]
bajo la medida neutral al riesgo. Para cada conjunto candidato de parámetros, los precios de opciones europeas bajo Heston se calculan mediante el método COS, utilizando la función característica del modelo. La función objetivo mide la diferencia entre los precios generados por Heston y los precios benchmark obtenidos a partir de la superficie de volatilidad implícita.\medskip

En paralelo, se construye la superficie de volatilidad local. Para ello, se transforma primero la superficie de volatilidad implícita en una superficie de precios call,
\[
V_c(K,T),
\]
y después se aplica la fórmula de Dupire en términos de precios:
\[
\sigma_{\mathrm{LV}}^2(K,T)
=
\frac{
\dfrac{\partial V_c(K,T)}{\partial T}
+
rK\dfrac{\partial V_c(K,T)}{\partial K}
}{
\dfrac{1}{2}K^2
\dfrac{\partial^2 V_c(K,T)}{\partial K^2}
}.
\]
Las derivadas parciales se aproximan numéricamente sobre la malla de strikes y vencimientos. Posteriormente, los valores de \(\sigma_{\mathrm{LV}}\) se interpolan para poder evaluarlos en puntos intermedios durante la simulación Monte Carlo.

\subsection{Estimación de la función de apalancamiento}

La simulación del modelo Heston--SLV requiere calcular en cada instante la función de apalancamiento
\[
\eta(K,t)
=
\frac{\sigma_{\mathrm{LV}}(K,t)}
{\sqrt{\mathbb{E}^{\mathbb{Q}}[v(t)\mid S(t)=K]}}.
\]
La esperanza condicional que aparece en el denominador no se conoce de forma cerrada, por lo que se aproxima con las trayectorias simuladas. En la implementación se utiliza un estimador por \textit{bins}: para cada instante \(t_j\), se ordenan las trayectorias según el valor del subyacente \(S_{i,j}\), se dividen en grupos con un número similar de trayectorias y se calcula la media de la varianza \(v_{i,j}\) dentro de cada grupo.\medskip

De esta forma, para cada trayectoria \(i\) y cada fecha \(t_j\), se obtiene una aproximación
\[
\widehat{E}_{i,j}
\approx
\mathbb{E}^{\mathbb{Q}}[v(t_j)\mid S(t_j)=S_{i,j}],
\]
y la función de apalancamiento aplicada en la trayectoria se calcula como
\begin{equation}
\label{eq:impl_slv_leverage}
\eta_{i,j}
=
\frac{\sigma_{\mathrm{LV}}(S_{i,j},t_j)}
{\sqrt{\max\{\widehat{E}_{i,j},\varepsilon\}}},
\end{equation}
donde \(\varepsilon>0\) es una constante pequeña que evita divisiones por valores demasiado próximos a cero.

\subsection{Simulación de trayectorias}

Una vez calibrados los parámetros de Heston y construida la superficie de volatilidad local, se simulan las trayectorias del par \((S(t),v(t))\). En la implementación se utiliza una malla temporal
\[
0=t_0<t_1<\cdots<t_n=T,
\qquad
\Delta t_j=t_{j+1}-t_j,
\]
que incluye todos los vencimientos necesarios para la valoración posterior.\medskip

El movimiento browniano de la varianza se genera a partir del movimiento browniano del subyacente mediante
\[
Z^v_{i,j}
=
\rho_{S,v} Z^S_{i,j}
+
\sqrt{1-\rho_{S,v}^2}\,Y_{i,j},
\]
donde \(Z^S_{i,j}\) e \(Y_{i,j}\) son normales estándar independientes. En el esquema de Euler, la varianza se actualiza mediante
\begin{equation}
\label{eq:impl_slv_var_euler}
v_{i,j+1}
=
\max\left\{
v_{i,j}
+
\kappa(\theta-\widetilde v_{i,j})\Delta t_j
+
\gamma\sqrt{\widetilde v_{i,j}}\sqrt{\Delta t_j}\,Z^v_{i,j},
0
\right\},
\end{equation}
donde \(\widetilde v_{i,j}=\max\{v_{i,j},\varepsilon\}\). El subyacente se simula de forma exponencial:
\begin{equation}
\label{eq:impl_slv_spot_euler}
S_{i,j+1}
=
S_{i,j}
\exp\left(
\left(r-\frac{1}{2}\eta_{i,j}^2\widetilde v_{i,j}\right)\Delta t_j
+
\eta_{i,j}\sqrt{\widetilde v_{i,j}}\sqrt{\Delta t_j}\,Z^S_{i,j}
\right).
\end{equation}

Además del esquema de Euler, el código permite utilizar un esquema de tipo \textit{Almost Exact Simulation} (AES) para la componente CIR de la varianza. En ese caso, la nueva varianza se obtiene a partir de la distribución de transición del proceso CIR, mientras que el subyacente se actualiza incorporando el término correlacionado con la variación de la varianza. Este esquema se incluye como alternativa numérica para reducir el sesgo de discretización, especialmente cuando se analizan vencimientos largos o pasos temporales más gruesos.\medskip

Tras cada paso temporal se aplica un ajuste de momentos al subyacente, de modo que la media empírica de las trayectorias coincida con el forward teórico:
\[
S_{i,j+1}
\leftarrow
S_{i,j+1}
\frac{S_0e^{rt_{j+1}}}{\frac{1}{N}\sum_{m=1}^N S_{m,j+1}}.
\]
Este ajuste ayuda a preservar la propiedad martingala del subyacente descontado en la simulación Monte Carlo.

\subsection{Valoración bajo Heston--SLV}

Una vez simuladas las trayectorias, la valoración de opciones europeas vainilla se realiza por promedio Monte Carlo. Para una call con strike \(K\) y vencimiento \(T=t_j\), se utiliza
\begin{equation}
\label{eq:impl_slv_vanilla}
\widehat{C}_{\mathrm{SLV}}(K,T)
=
e^{-rT}
\frac{1}{N}
\sum_{i=1}^{N}
\max\{S_{i,j}-K,0\}.
\end{equation}
El caso put se obtiene sustituyendo el payoff por \(\max\{K-S_{i,j},0\}\).\medskip

Para las opciones \textit{forward-start}, la implementación utiliza la versión normalizada por el valor del subyacente en la fecha inicial del periodo forward. Si \(0<T_1<T_2\), y ambos tiempos están incluidos en la simulación, se calcula el cociente
\[
R_i(T_1,T_2)
=
\frac{S_i(T_2)}{S_i(T_1)}.
\]
Así, para una call \textit{forward-start} con strike relativo \(K_{\mathrm{rel}}\), el precio se estima mediante
\begin{equation}
\label{eq:impl_slv_forward}
\widehat{C}_{\mathrm{SLV}}^{FS}(K_{\mathrm{rel}},T_1,T_2)
=
e^{-rT_2}
\frac{1}{N}
\sum_{i=1}^{N}
\max\left\{
\frac{S_i(T_2)}{S_i(T_1)}
-
K_{\mathrm{rel}},
0
\right\}.
\end{equation}
Esta convención es la que se utilizará también para comparar con la metodología MRE.

\section{Implementación de Entropía Relativa Mínima}

\subsection{Simulación de la medida previa}

La metodología MRE parte de una medida previa discreta. En la implementación, dicha medida se obtiene simulando trayectorias de un movimiento browniano geométrico bajo la medida neutral al riesgo:
\begin{equation}
\label{eq:impl_mre_prior}
dS(t)=(r-q)S(t)\,dt+\sigma(t)S(t)\,dW(t).
\end{equation}
La volatilidad puede tomarse constante o puede construirse a partir de una estructura temporal implícita at-the-money-forward extraída de los benchmarks. En el segundo caso, se infiere una varianza total para cada vencimiento y se obtienen volatilidades por intervalo temporal. Así, el prior mantiene una dinámica sencilla, pero recoge parcialmente el nivel temporal de volatilidad observado en los datos.\medskip

Las trayectorias se simulan en las fechas necesarias para calcular los instrumentos de calibración. Si se quieren valorar productos dependientes de dos fechas, como opciones \textit{forward-start}, esas fechas deben estar incluidas en la matriz de trayectorias o añadirse a la malla de simulación.

\subsection{Construcción de la matriz de restricciones}

Para cada trayectoria \(i\) y cada instrumento de calibración \(j\), se calcula el payoff descontado
\[
g_{ij}.
\]
Con estos valores se construye la matriz
\[
G=(g_{ij})_{1\leq i\leq N,\;1\leq j\leq m}.
\]
Los precios objetivo se agrupan en un vector
\[
C=(C_1,\ldots,C_m).
\]
En la implementación se consideran tres tipos de restricciones:
\[
\sum_{i=1}^N q_i g_{ij}=C_j.
\]

En primer lugar, se incluyen restricciones de opciones vainilla, cuyo objetivo es reproducir los precios benchmark. En segundo lugar, se pueden añadir restricciones forward, para controlar la media descontada del subyacente en cada vencimiento. Finalmente, se incorporan restricciones sintéticas de martingala por ventanas, siguiendo la idea de que la calibración a opciones de distintos vencimientos no garantiza por sí sola que el proceso ponderado preserve la propiedad martingala.\medskip

Estas restricciones auxiliares no representan opciones cotizadas, sino condiciones de consistencia financiera. Su papel es especialmente importante cuando la medida posterior se va a utilizar para valorar derivados no incluidos en la calibración.

\subsection{Resolución del problema dual}

Sea \(p_i\) el peso previo de la trayectoria \(i\). La calibración busca pesos posteriores \(q_i\) que reproduzcan las restricciones impuestas y que se alejen lo mínimo posible de la medida previa. En la implementación se utiliza la forma exponencial de la solución:
\begin{equation}
\label{eq:impl_mre_weights}
q_i(\lambda)
=
\frac{
p_i\exp\left(\sum_{j=1}^m \lambda_j g_{ij}\right)
}{
\sum_{\ell=1}^N
p_\ell\exp\left(\sum_{j=1}^m \lambda_j g_{\ell j}\right)
}.
\end{equation}
Por tanto, el problema se reduce a encontrar los multiplicadores \(\lambda\) que minimizan la función dual. El código contempla tanto el ajuste exacto como la formulación de mínimos cuadrados ponderados. Esta segunda formulación permite relajar restricciones mediante parámetros de tolerancia y resulta más estable cuando el número de restricciones es elevado o cuando algunas restricciones son casi redundantes.\medskip

La minimización se realiza mediante un método de Newton basado en el gradiente y la Hessiana de la función dual. Para mejorar la estabilidad numérica se evalúa la normalización de los pesos de forma estable, se amortigua el paso de Newton cuando es necesario y se añade una pequeña regularización diagonal al sistema lineal. Si el sistema está mal condicionado, se utiliza una resolución por mínimos cuadrados.

\subsection{Valoración bajo la medida posterior}

Una vez obtenidos los pesos calibrados \(q_i\), la medida posterior queda determinada sobre el conjunto de trayectorias simuladas. Para una call europea con vencimiento \(T=t_j\), el estimador bajo MRE es
\begin{equation}
\label{eq:impl_mre_vanilla}
\widehat{C}_{\mathrm{MRE}}(K,T)
=
\sum_{i=1}^N
q_i
e^{-rT}
\max\{S_{i,j}-K,0\}.
\end{equation}
La diferencia con el estimador Monte Carlo usual no está en el payoff, sino en la medida utilizada: en lugar de promediar todas las trayectorias con peso \(1/N\), se ponderan mediante los pesos posteriores \(q_i\).\medskip

Para una opción \textit{forward-start} normalizada, el mismo principio conduce a
\begin{equation}
\label{eq:impl_mre_forward}
\widehat{C}_{\mathrm{MRE}}^{FS}(K_{\mathrm{rel}},T_1,T_2)
=
\sum_{i=1}^N
q_i
e^{-rT_2}
\max\left\{
\frac{S_i(T_2)}{S_i(T_1)}
-
K_{\mathrm{rel}},
0
\right\}.
\end{equation}
Esta expresión permite valorar el mismo producto que en \eqref{eq:impl_slv_forward}, pero utilizando la medida posterior calibrada por entropía relativa.

\section{Diagnósticos numéricos}

Antes de comparar ambos métodos, se calculan varios diagnósticos para comprobar que las implementaciones son coherentes. En el caso MRE, se verifica que los pesos posteriores suman uno, que son no negativos y que los errores de calibración frente a las restricciones impuestas son pequeños. También se registran por separado los errores asociados a opciones, forwards y restricciones de martingala.\medskip

En el caso Heston--SLV, se comprueba la estabilidad de los precios Monte Carlo frente al número de pasos temporales y al esquema de simulación utilizado. La comparación entre Euler y AES debe interpretarse como un diagnóstico de discretización del método SLV, no como un resultado financiero central del trabajo. Por tanto, una tabla de este tipo puede incluirse en el capítulo de resultados como análisis de robustez, o bien en un apéndice si se quiere mantener la comparación principal centrada en SLV frente a MRE.\medskip

Asimismo, el código permite comparar la sonrisa \textit{forward-start} generada por Heston puro, volatilidad local pura y Heston--SLV. Esta gráfica es útil para comprobar que el modelo SLV combina el ajuste inicial de la volatilidad local con una dinámica estocástica más rica. Sin embargo, no constituye la comparación principal del trabajo, ya que el objetivo final es contrastar los precios obtenidos mediante SLV con los obtenidos mediante Entropía Relativa Mínima. Por este motivo, dicha gráfica se utilizará únicamente como apoyo o diagnóstico del bloque SLV, mientras que las tablas principales del capítulo de resultados se centrarán en comparar SLV y MRE sobre los mismos instrumentos.
